\scanpage{86}
% archive foliation: 85
Il y a une telle structure canonique dans le cas de l'exemple.
Ou aussi : quand peut-on enrichir la composition des morphismes de $\unsure{\mathcal P|P}$
\[\deleted{\operatorname{Hom}_\unsure{\mathcal P|P}}\unsure{\mathcal P|P}(X,Y)\times \unsure{\mathcal P|P}(Y,Z)\mathbin{-}\unsure{\mathcal P|P}(X,Z)\]
($k$-linéaire en le deuxième argument) en une application linéaire
\[\unsure{\mathcal P|P}(Y,Z)\longrightarrow(\text{Appl. }k\text{-polynomiales})(\unsure{\mathcal P|P}(X,Y),\unsure{\mathcal P|P}(X,Z))\]
(ceci aussi est acquis dans l'exemple \unsure{principal} envisagé). La construction doit se ramener pour l'essentiel à définir de telles structures suppl. sur les foncteurs $\Gamma^n$, i.e. pour tt $X,Y$
\[f\longmapsto\Gamma^p(f):\operatorname{Hom}_1(X,Y)\mathbin{-}\operatorname{Hom}_1(\Gamma^pX,\Gamma^pY)\]
doit être défini comme ``application pol. de degré $p$'' entre les $k$-modules $\operatorname{Hom}_1(X,Y)$ et $\operatorname{Hom}_1(\Gamma^{\unsure{r|p}}X,\Gamma^{\unsure{r|p}}Y)$, et non seulement comme application ensembliste.

